\documentclass[10pt,a4paper]{article}

% Die Loesungsfassung unterscheidet sich ausschliesslich durch \solutiontrue.
\usepackage[T1]{fontenc}
\usepackage[utf8]{inputenc}
\usepackage[ngerman]{babel}
\usepackage[a4paper,margin=14mm]{geometry}
\usepackage{lmodern}
\usepackage[table]{xcolor}
\usepackage{tabularx}
\usepackage{microtype}
\usepackage{amssymb}
\usepackage{tikz}

\newif\ifsolution
\solutionfalse

\pagestyle{empty}
\setlength{\parindent}{0pt}
\setlength{\parskip}{0pt}
\renewcommand{\familydefault}{\sfdefault}

\definecolor{Navy}{HTML}{17324D}
\definecolor{Blue}{HTML}{3B6EA5}
\definecolor{Gold}{HTML}{F3B33D}
\definecolor{Ink}{HTML}{1E2935}
\definecolor{Muted}{HTML}{617080}
\definecolor{Line}{HTML}{D8E1E8}
\color{Ink}

\newcommand{\kopf}{%
  \colorbox{Navy}{\parbox{\dimexpr\linewidth-2\fboxsep\relax}{%
    \vspace{1.3mm}\color{white}
    \begin{tabularx}{\linewidth}{@{}Xr@{}}
      {\small\bfseries PHYSIK · KLASSE 8} &
      \colorbox{Gold}{\color{Navy}\small\bfseries\strut\ifsolution LÖSUNGEN\else TEST\fi}
    \end{tabularx}
    \vspace{0.9mm}
    {\fontsize{16}{19}\selectfont\bfseries Elektrischer Strom · Spannung messen}\par
    \vspace{0.3mm}{\small Test zur Abfrage · Förderbandmodell und Spannungsmessung}\vspace{1.1mm}
  }}\par\vspace{2mm}}
\newcommand{\frage}[2]{%
  \vspace{0.7mm}{\color{Blue}\bfseries\normalsize #1 · #2}\par
  \vspace{0.2mm}{\color{Blue}\rule{\linewidth}{0.55pt}}\par\vspace{0.4mm}}
\newcommand{\wahl}[2]{%
  \noindent\ifsolution\ifnum#1=1 $\boxtimes$\else $\square$\fi\else $\square$\fi\enspace #2\par}
\newcommand{\antwortfeld}[2]{%
  \vspace{0.7mm}%
  \begin{tikzpicture}[x=1mm,y=1mm]
    \path[use as bounding box] (0,0) rectangle (181.5,-7*#1);
    \ifsolution
      \node[anchor=north west,text width=179mm,inner sep=0pt,align=left,
        font=\small,text=Blue] at (0,-1) {#2};
    \else
      \foreach \i in {1,...,#1}{\draw[Line,line width=.55pt] (0,-7*\i) -- (181.5,-7*\i);}
    \fi
  \end{tikzpicture}\par}
\newcommand{\spannungsschaltung}[2]{%
  \begin{scope}[shift={(#1,#2)}]
    \draw[Blue,line width=.85pt]
      (0,0)--(31,0) (34,0)--(72,0)--(72,-12)--(39.5,-12)
      (0,0)--(0,-12)--(32.5,-12);
    \draw[Blue,line width=.85pt] (31,-3.5)--(31,3.5) (34,-2)--(34,2);
    \node[font=\scriptsize,text=Blue] at (28,4) {$+$};
    \node[font=\scriptsize,text=Blue] at (37,4) {$-$};
    \draw[Blue,line width=.85pt,fill=white] (36,-12) circle (3.5);
    \draw[Blue,line width=.7pt] (33.5,-14.5)--(38.5,-9.5) (33.5,-9.5)--(38.5,-14.5);
    \draw[Blue,line width=.85pt]
      (16,-12)--(16,-24)--(32.5,-24) (39.5,-24)--(56,-24)--(56,-12);
    \draw[Blue,line width=.85pt,fill=white] (36,-24) circle (3.5);
    \node[font=\small,text=Blue] at (36,-24) {V};
    \foreach \x in {16,56}{\fill[Blue] (\x,-12) circle (1);}
    \node[anchor=west,font=\small,text=Blue] at (78,0) {Batterie};
    \node[anchor=west,font=\small,text=Blue] at (78,-12) {Glühbirne};
  \end{scope}}
\newcommand{\karofeld}[2]{%
  \vspace{0.5mm}%
  \begin{tikzpicture}[x=1mm,y=1mm]
    \path[use as bounding box] (0,0) rectangle (181.5,-#1);
    \draw[step=5mm,Line!65,line width=.25pt] (0,-#1) grid (181.5,0);
    \draw[Line,line width=.55pt] (0,0) rectangle (181.5,-#1);
    \ifsolution #2\fi
  \end{tikzpicture}\par}
\newcommand{\umrechnungen}[6]{%
  \begin{tikzpicture}[x=1mm,y=1mm]
    \path[use as bounding box] (0,0) rectangle (181.5,-8);
    \node[anchor=west,inner sep=0pt,text=\ifsolution Blue\else Ink\fi] at (0,-4)
      {a) $#1=\ifsolution #2\else\rule{24mm}{.4pt}\fi\,\mathrm{#3}$};
    \node[anchor=west,inner sep=0pt,text=\ifsolution Blue\else Ink\fi] at (91,-4)
      {b) $#4=\ifsolution #5\else\rule{24mm}{.4pt}\fi\,\mathrm{#6}$};
  \end{tikzpicture}\par}
\newcommand{\fuss}{%
  \vfill{\color{Line}\rule{\linewidth}{0.5pt}}\par\vspace{1mm}
  \begin{tabularx}{\linewidth}{@{}Xr@{}}
    {\small\color{Muted}Klasse 8 · Physik · Elektrischer Strom} &
    {\small\color{Muted}\ifsolution Lösungen\else Test\fi{} · Seite 1 von 1}
  \end{tabularx}}

\begin{document}
\kopf
\textbf{Untersuche} das Förderbandmodell als Modell für den elektrischen Stromkreis.\par

\frage{1}{Ampere und Milliampere}
\textbf{Rechne} die beiden Stromstärken in die jeweils angegebene Einheit um.\par
\umrechnungen{240\,\mathrm{mA}}{0{,}24}{A}{0{,}06\,\mathrm{A}}{60}{mA}

\frage{2}{Kreuze alle richtigen Zuordnungen im Förderbandmodell an.}
\wahl{1}{Im Modell stellt das Förderband die Batterie dar.}
\wahl{0}{Im Modell stellt das Drehkreuz die Batterie dar.}
\wahl{0}{Im Modell stellt der Höhenunterschied die Stromstärke dar.}
\wahl{1}{Im Modell stellt der Höhenunterschied die elektrische Spannung dar.}

\frage{3}{Kreuze alle richtigen Aussagen zu Kugeln und Drehkreuz an.}
\wahl{0}{Im Modell werden die Kugeln am Drehkreuz verbraucht.}
\wahl{1}{Im Modell stellt das Drehkreuz die Glühbirne dar.}
\wahl{1}{Im Modell kehren die Kugeln nach dem Drehkreuz zum Förderband zurück.}
\wahl{0}{Im Modell erzeugt das Förderband ständig neue Kugeln.}

\frage{4}{Kreuze alle richtigen Aussagen zum Spannungsmessgerät an.}
\wahl{1}{Die elektrische Spannung wird mit einem Voltmeter gemessen.}
\wahl{0}{Das Schaltzeichen des Voltmeters ist ein Kreis mit einem A.}
\wahl{1}{Das Schaltzeichen des Voltmeters ist ein Kreis mit einem V.}
\wahl{0}{Ein Amperemeter misst die elektrische Spannung.}

\frage{5}{Kreuze alle richtigen Aussagen zum Anschluss des Voltmeters an.}
\wahl{1}{Für die Spannung an der Lampe wird das Voltmeter parallel zur Lampe angeschlossen.}
\wahl{0}{Für die Spannung an der Lampe wird das Voltmeter in Reihe zur Lampe eingebaut.}
\wahl{0}{Ein einziger Anschluss des Voltmeters an die Lampe genügt zur Spannungsmessung.}
\wahl{1}{Für die Spannung an der Batterie wird das Voltmeter mit beiden Batteriepolen verbunden.}

\frage{6}{Die Kugel am Drehkreuz}
\textbf{Beschreibe} im Förderbandmodell die Wirkung einer Kugel am Drehkreuz und ihren weiteren Weg. \textbf{Erkläre} den Bezug zu der Situation im elektrischen Stromkreis.\par
\antwortfeld{4}{Im Modell treibt die Kugel das Drehkreuz an und rollt zum Förderband zurück. Die Kugeln stehen für Elektronen, das Drehkreuz für die Glühbirne. Die Elektronen geben Energie an die Glühbirne ab; diese wandelt elektrische Energie in Licht und Wärme um. Die Elektronen werden dabei nicht verbraucht und bewegen sich im geschlossenen Stromkreis weiter.}

\frage{7}{Die Spannung an der Glühbirne messen}
\textbf{Skizziere} einen geschlossenen Stromkreis aus Batterie und Lampe. Es soll die Spannung gemessen werden, die an der Glühbirne abfällt. \textbf{Ergänze} ein entsprechendes Messgerät.\par
\karofeld{32}{\begin{scope}[scale=.85]\spannungsschaltung{47}{-6}\end{scope}}

\fuss
\end{document}
